\section{C6 under loading, preparation and native stepping}\label{sec:preparation}
The stateful model earns its account from the component dynamics. It does not remove the temporal-adjoint result of Section~\ref{sec:observation}. We now identify precisely the retained observation and the effect of preparing independent line states.

\subsection{Initial-value and compatible-history formulas}
For nonnegative-time histories define
\begin{equation}\label{eq:delayplus}
(D_\theta^+h)(s,X,t)=\ii_{t\ge\theta_X(s)}h(s,X,t-\theta_X(s)),\qquad
c_{\rm init}(s,X,t)=\ii_{0\le t<\theta_X(s)}u_{+,0}(s,X,\theta_X(s)-t).
\end{equation}
The startup term is L2 in time by Lemma~\ref{lem:transport}. Substituting the actual near boundary law into~\eqref{eq:cstartup} gives
$c=\sqrt2\kappa D_\theta^+JQ-D_\theta^+b+c_{\rm init}$.
Applying C and the receiving ODE proves the exact identity
\begin{equation}\label{eq:loadedC6}
\boxed{\begin{aligned}
\beta_G\dot\zeta&=2\kappa\KGz Q-\sqrt2CD_\theta^+b
+\sqrt2C c_{\rm init}-CC^*\zeta,\\
\KGz&=CD_\theta^+J.
\end{aligned}}
\end{equation}
This holds in locally L2 time with values in Y for every finite-energy preparation. The startup contribution has a plus sign, and the two $\sqrt2$ factors produce the direct coefficient $2\kappa$. Here $\KGz$ is C6 with a specified zero extension of Q as an algebraic decomposition. That extension does not infer that an arbitrary physical past was zero. Independent forward profiles carry separate preparation and need not coincide with a past RX trajectory.

\begin{proposition}[Compatible complete histories]\label{prop:completehistory}
Suppose the required negative-time histories are supplied, with
\begin{equation}\label{eq:historycompatibility}
u_{+,0}(s,X,\xi)=a(s,X,-\xi),\qquad
a+b=\sqrt2\kappa JQ
\end{equation}
on those times as well as forward in time. Then
\begin{equation}\label{eq:completeC6}
\boxed{\beta_G\dot\zeta=2\kappa\KG Q-\sqrt2CD_\theta b-CC^*\zeta.}
\end{equation}
\end{proposition}
\begin{proof}
For $t<\theta$, the compatible startup expression is $a(s,X,t-\theta)$; for $t\ge\theta$, the same expression comes from the boundary. Thus $c=D_\theta a$. Substitute the negative-time and forward near law and use~\eqref{eq:factorization}. Suitable L2 histories on the needed finite interval justify each term by time translation.
\end{proof}

A Q history alone does not imply~\eqref{eq:historycompatibility}, which also involves a and b. Supplied histories satisfying an identity are not automatically histories generated by the full new model; no backward-existence theorem has been used. After $t\ge\theta_+$ the startup term in~\eqref{eq:loadedC6} vanishes a.e. and all its delayed arguments are nonnegative.

Both formulas are causal. Their $C^*$ is a same-time spatial adjoint, not the finite-window temporal $\KG^*$. C6 remains an exact direct component of the actual Q history, but neither $\zeta$ nor its full loaded drive is generally bare C6. Return and loading, and for arbitrary preparation startup, remain explicit. Nonnegative $E_{\rm tr}(0)$ is a quadratic stock of new states; it is not the signed bilinear $W_{\rm pre}$ of~\eqref{eq:testpre} under a new name.

\subsection{The unchanged native map and finite hybrid}\label{sec:native}
Retain the exact original tuple~\cite{RXhybrid}
\begin{equation}\label{eq:nativetuple}
\epsilon=\frac1{10},\quad g=\frac1{60},\quad \lambda_N=\frac1{300},\quad
k=\left(\frac{999}{1000},1,\frac{1001}{1000}\right).
\end{equation}
For $\Omega\in\CC^3$, all components use the same original state for the simultaneous prestage and the same prestage for synchronization:
\begin{align}
\widetilde\Omega_i&=\Omega_i+\epsilon(k_i-|\Omega_i|^2)\Omega_i
+g\sum_{j\ne i}(\Omega_j-\Omega_i),\label{eq:prestage}\\
\delta_i&=\lambda_N\sum_{j\ne i}\sin\!\left(3[\operatorname{Arg}_0\widetilde\Omega_j-\operatorname{Arg}_0\widetilde\Omega_i]\right),
&\Fn(\Omega)_i&=e^{i\delta_i}\widetilde\Omega_i.\label{eq:nativemap}
\end{align}
The convention is $\operatorname{Arg}_0(0)=0$; exact zero prestage components remain zero after rotation. No normalization, coefficient change or continuous-gradient substitute is made.

At a native jump apply this map only to x and leave Q,z,$u_+,u_-,\zeta$ unchanged. Then perform the new exchange stage for a fixed symbolic positive dwell. This is F-first/exchange-second, not a numerical integration scheme.

\begin{proposition}[Finite native-step account]\label{prop:hybrid}
Every finite native/exchange iterate is defined for finite initial energy. For a finite window with post-jump endpoint values,
\begin{equation}\label{eq:hybridbalance}
\boxed{H_{\rm total}(t_b)-H_{\rm total}(t_a)
=\sum_{t_a<t_j\le t_b}\frac{\alpha_N}{2}
\left[I(\Fn\Omega(t_j^-))-I(\Omega(t_j^-))\right]=B_N.}
\end{equation}
\end{proposition}
\begin{proof}
Only native storage changes at each jump; its increment is the summand. Exchange conserves total stock by Theorem~\ref{thm:global}. Summation with the source endpoint convention proves~\eqref{eq:hybridbalance}. A jump at the left endpoint is already in its post-jump value, while one at the right is counted. Finite native inputs have finite polynomial prestages and defined finite images under the zero convention. The global exchange theorem supplies the next interval. Induction gives every finite iterate, without a uniform infinite-step bound.
\end{proof}

The baseline increment is not additional port work or automatically a specified external physical source. No extra impulses are introduced. Nor does finite existence imply continuity across singular native prestages, preservation of a coherent domain, an increasing-step-count limit or recovery after repeated preloaded native updates. Those qualifications remain as in the accepted cadence work~\cite{CAD}.

\subsection{True zero-coupling recovery and D01}
At $\kappa=0$, $\FF=0$ for every incoming b, so the original RX subsystem is recovered for \emph{every} added-state preparation. The isolated new subsystem has $a=-b$, hence $P_N=0$, and can evolve using its own conserved $E_{\rm tr}+H_G$. Empty added states stay empty by uniqueness.

\begin{remark}[Stationary nonempty preparation; D01]\label{rem:D01}
Empty added states are sufficient, not necessary, for inactivity at zero coupling. Choose nonzero $h\in L^2(G)$ and a unit normal n to the carrier plane, and prepare
\begin{equation}\label{eq:inactivekernel}
\kappa=0,\qquad u_{+,0}=u_{-,0}=0,\qquad \zeta_0(X)=(0,h(X)n)\ne0.
\end{equation}
Since $C^*\zeta_0=CC^*\zeta_0=0$, all boundary waves remain zero and $\zeta(t)=\zeta_0$ has positive constant stock. The old RX preparation can be arbitrary and decoupled. This is D01's correction to an earlier work-order necessity sentence; the candidate equations and analytical theorem are unchanged~\cite{Treview}.
\end{remark}

At positive coupling the force changes Q and hence the coupled trajectory. Keeping the native and old-z equations in the same algebraic form does not automatically transfer their earlier capacity, credit, convergence or receiving/recovery theorems. Their actual hypotheses and invariants must be checked again for such a use.
